% Standalone build: open this file in VS Code and run the build, or
%   latexmk -cd en/lessons/01-using-this-template.tex
\documentclass[../main.tex]{subfiles}
\begin{document}

\chapter{Using this template}
\lessoninfo{
  video={Lesson 1 (Introduction)},
  textbook={H\&P \cite{hennessy2017} §1.1--1.3},
  keywords={sidenotes, index, bilingual build},
}

This chapter is a showcase of every macro the class provides.  Delete it
when you start a real course, but keep \texttt{02-lesson-skeleton.tex} as a
starting point for new lessons.

\section{Margin notes}

The right-hand column is for anything that supports the main text without
interrupting it.  Numbered remarks use \verb|\sidenote|\sidenote{Like this
one.  Numbers restart at every chapter.} and unnumbered remarks use
\verb|\margin|.\margin{An unnumbered margin note.}

Three labelled variants make the \emph{kind} of side content visible at a
glance:
\source{Amdahl's original paper~\cite{amdahl1967}; H\&P §1.9.}
\tip{Use \texttt{\string\tip} for exam-relevant advice or mnemonics.}
\caution{Use \texttt{\string\caution} for common misconceptions.}
\begin{itemize}
  \item \verb|\source{...}| --- where the material comes from (lecture, textbook section, paper);
  \item \verb|\tip{...}| --- practical advice;
  \item \verb|\caution{...}| --- pitfalls and misconceptions.
\end{itemize}

\section{Terms and the index}

A \term{pipeline}[パイプライン] is introduced with \verb|\term|, which
emphasizes the word, adds it to the index, and optionally prints a gloss in
the margin (here: the Japanese term).  Later mentions of a
\term*{pipeline} use \verb|\term*| so the index does not fill up with
repeated page numbers.  Terms can be nested: \term{speedup}[高速化率] and
\term{Amdahl's law}[Amdahl の法則] both go to the index.

\begin{definition}[Speedup]
  For a change that makes a fraction $f$ of the execution time $s$ times
  faster, the overall speedup is
  \begin{equation}
    S = \frac{1}{(1-f) + f/s}.
    \label{eq:amdahl}
  \end{equation}
\end{definition}

\begin{example}
  If $f = 0.9$ and $s \to \infty$, \cref{eq:amdahl} gives $S = 10$: the
  untouched \SI{10}{\percent} bounds the achievable speedup.
\end{example}

\begin{note}
  \verb|definition|, \verb|example| and \verb|note| share one counter per
  chapter, so they can be cross-referenced with \verb|\cref|.
\end{note}

\begin{keypoint}[Amdahl's law]
  Optimize the common case; the part you do not touch limits the gain.
  Margin macros work inside this box too.\margin{From inside a key point.}
\end{keypoint}

\section{Figures, tables and code}

\begin{marginfigure}
  \begin{tikzpicture}[scale=0.6]
    \draw[->] (0,0) -- (4,0) node[right,font=\tiny] {$f$};
    \draw[->] (0,0) -- (0,3) node[above,font=\tiny] {$S$};
    \draw[thick,ln-accent] plot[domain=0:3.6,samples=40] (\x,{1/(1-\x/4+0.001)*0.7});
  \end{tikzpicture}
  \caption{A figure in the margin (\texttt{marginfigure}).}
  \label{fig:margin}
\end{marginfigure}

Small figures go in the margin (\cref{fig:margin}); normal figures use the
standard \verb|figure| environment; wide figures use \verb|widefigure|,
which spans the body and the margin column (\cref{fig:wide}).  Figures shared
between the two editions live in \texttt{figures/} and switch their labels
with \verb|\iflangja|.

\begin{figure}[htbp]
  \centering
  % Shared TikZ figure.  Labels switch with the document language.
% Usage: % Shared TikZ figure.  Labels switch with the document language.
% Usage: % Shared TikZ figure.  Labels switch with the document language.
% Usage: \input{figures/pipeline-stages}
\begin{tikzpicture}[
  stage/.style={draw, rounded corners=1pt, minimum width=1.6cm, minimum height=0.8cm,
                font=\sffamily\small, fill=ln-box},
  >={Stealth[length=2mm]}, node distance=4mm,
]
  \node[stage] (if)  {IF};
  \node[stage, right=of if]  (id)  {ID};
  \node[stage, right=of id]  (ex)  {EX};
  \node[stage, right=of ex]  (mem) {MEM};
  \node[stage, right=of mem] (wb)  {WB};
  \draw[->] (if) -- (id);
  \draw[->] (id) -- (ex);
  \draw[->] (ex) -- (mem);
  \draw[->] (mem) -- (wb);
  \node[below=1mm of if,  font=\footnotesize\color{ln-muted}] {\iflangja{命令フェッチ}{fetch}};
  \node[below=1mm of id,  font=\footnotesize\color{ln-muted}] {\iflangja{デコード}{decode}};
  \node[below=1mm of ex,  font=\footnotesize\color{ln-muted}] {\iflangja{実行}{execute}};
  \node[below=1mm of mem, font=\footnotesize\color{ln-muted}] {\iflangja{メモリ}{memory}};
  \node[below=1mm of wb,  font=\footnotesize\color{ln-muted}] {\iflangja{書き戻し}{write back}};
\end{tikzpicture}

\begin{tikzpicture}[
  stage/.style={draw, rounded corners=1pt, minimum width=1.6cm, minimum height=0.8cm,
                font=\sffamily\small, fill=ln-box},
  >={Stealth[length=2mm]}, node distance=4mm,
]
  \node[stage] (if)  {IF};
  \node[stage, right=of if]  (id)  {ID};
  \node[stage, right=of id]  (ex)  {EX};
  \node[stage, right=of ex]  (mem) {MEM};
  \node[stage, right=of mem] (wb)  {WB};
  \draw[->] (if) -- (id);
  \draw[->] (id) -- (ex);
  \draw[->] (ex) -- (mem);
  \draw[->] (mem) -- (wb);
  \node[below=1mm of if,  font=\footnotesize\color{ln-muted}] {\iflangja{命令フェッチ}{fetch}};
  \node[below=1mm of id,  font=\footnotesize\color{ln-muted}] {\iflangja{デコード}{decode}};
  \node[below=1mm of ex,  font=\footnotesize\color{ln-muted}] {\iflangja{実行}{execute}};
  \node[below=1mm of mem, font=\footnotesize\color{ln-muted}] {\iflangja{メモリ}{memory}};
  \node[below=1mm of wb,  font=\footnotesize\color{ln-muted}] {\iflangja{書き戻し}{write back}};
\end{tikzpicture}

\begin{tikzpicture}[
  stage/.style={draw, rounded corners=1pt, minimum width=1.6cm, minimum height=0.8cm,
                font=\sffamily\small, fill=ln-box},
  >={Stealth[length=2mm]}, node distance=4mm,
]
  \node[stage] (if)  {IF};
  \node[stage, right=of if]  (id)  {ID};
  \node[stage, right=of id]  (ex)  {EX};
  \node[stage, right=of ex]  (mem) {MEM};
  \node[stage, right=of mem] (wb)  {WB};
  \draw[->] (if) -- (id);
  \draw[->] (id) -- (ex);
  \draw[->] (ex) -- (mem);
  \draw[->] (mem) -- (wb);
  \node[below=1mm of if,  font=\footnotesize\color{ln-muted}] {\iflangja{命令フェッチ}{fetch}};
  \node[below=1mm of id,  font=\footnotesize\color{ln-muted}] {\iflangja{デコード}{decode}};
  \node[below=1mm of ex,  font=\footnotesize\color{ln-muted}] {\iflangja{実行}{execute}};
  \node[below=1mm of mem, font=\footnotesize\color{ln-muted}] {\iflangja{メモリ}{memory}};
  \node[below=1mm of wb,  font=\footnotesize\color{ln-muted}] {\iflangja{書き戻し}{write back}};
\end{tikzpicture}

  \caption{A classic five-stage pipeline (shared TikZ source).}
  \label{fig:pipeline}
\end{figure}

\begin{widefigure}[htbp]
  \begin{tabular}{@{}lccccccc@{}}
    \toprule
    Cycle & 1 & 2 & 3 & 4 & 5 & 6 & 7 \\
    \midrule
    \texttt{add}  & IF & ID & EX & MEM & WB &    &    \\
    \texttt{sub}  &    & IF & ID & EX  & MEM & WB &   \\
    \texttt{lw}   &    &    & IF & ID  & EX  & MEM & WB \\
    \bottomrule
  \end{tabular}
  \caption{A wide float spanning body and margin (\texttt{widefigure}).}
  \label{fig:wide}
\end{widefigure}

Code uses \verb|listings|:
\begin{lstlisting}[language={[x86masm]Assembler}, caption={A loop in assembly.}]
loop:   add   r1, r1, r2      ; accumulate
        addi  r3, r3, -1
        bne   r3, r0, loop    ; branch if not done
\end{lstlisting}

Pseudocode uses \verb|algorithmicx|:
\begin{algorithm}[htbp]
  \caption{Two-bit saturating counter update}
  \begin{algorithmic}[1]
    \If{taken}
      \State $c \gets \min(c+1, 3)$
    \Else
      \State $c \gets \max(c-1, 0)$
    \EndIf
  \end{algorithmic}
\end{algorithm}

\section{Bilingual writing}

Both editions share \texttt{format/}, \texttt{figures/} and
\texttt{references.bib}; only the prose under \texttt{en/} and \texttt{ja/}
differs.  Inside shared files use \verb|\iflangja{japanese}{english}| or the
shorthands \verb|\ja{...}| and \verb|\en{...}|.

\end{document}
